\documentclass[aspectratio=169,10pt]{beamer}
\usetheme[numbering=fraction,progressbar=frametitle]{metropolis}
\usepackage{fontspec}
\setsansfont{Latin Modern Sans}
\usepackage{amsmath,amssymb,mathtools,graphicx,booktabs}
\usepackage{polyglossia}
\setdefaultlanguage{french}
\setbeamertemplate{navigation symbols}{}
\setbeamerfont{frametitle}{size=\large}
\setbeamerfont{normal text}{size=\normalsize}
\setlength{\parskip}{.45em}
\hypersetup{pdfauthor={},pdfsubject={Lecture slides},pdfkeywords={},pdftitle={Fonction logarithme décimal}}
\providecommand{\N}{\mathbb{N}}
\providecommand{\Z}{\mathbb{Z}}
\providecommand{\D}{\mathbb{D}}
\providecommand{\Q}{\mathbb{Q}}
\providecommand{\R}{\mathbb{R}}
\providecommand{\C}{\mathbb{C}}
\providecommand{\vect}[1]{\overrightarrow{#1}}
\providecommand{\vi}{\vec{\imath}}
\providecommand{\vj}{\vec{\jmath}}
\providecommand{\vk}{\vec{k}}
\providecommand{\norm}[1]{\left\lVert #1\right\rVert}
\providecommand{\vc}[2]{\begin{pmatrix}#1\\#2\end{pmatrix}}
\providecommand{\deter}[2]{\det\left(#1,#2\right)}
\providecommand{\abs}[1]{\left\lvert #1\right\rvert}
\providecommand{\dd}{\,\mathrm{d}}
\providecommand{\e}{\mathrm{e}}
\providecommand{\Card}{\operatorname{Card}}
\providecommand{\E}{\operatorname{E}}
\providecommand{\V}{\operatorname{V}}
\providecommand{\degre}{^{\circ}}
\providecommand{\un}[1]{\,\mathrm{#1}}
\providecommand{\sci}[2]{#1\times10^{#2}}
\providecommand{\pgcd}{\operatorname{PGCD}}
\providecommand{\ppcm}{\operatorname{PPCM}}
\title{Fonction logarithme décimal}
\subtitle{Cours et exemples}
\date{}
\author{}
\begin{document}
\begin{frame}[plain]
\vfill
{\usebeamerfont{title}\inserttitle\par}
\medskip
{\usebeamerfont{subtitle}\insertsubtitle\par}
\medskip
{\usebeamercolor[fg]{progress bar}\rule{\textwidth}{.5pt}}
\vfill
\end{frame}
\begin{frame}{Cours de référence}
Pour tout \(b>0\), \(\log(b)\) est l'unique réel \(x\) tel que \(10^x=b\). Ainsi \(\log(10^x)=x\) et, pour \(b>0\), \(10^{\log(b)}=b\). La fonction \(\log\) est définie sur \(]0,+\infty[\) et elle est strictement croissante.
\end{frame}
\begin{frame}{Cours de référence}
\(\log(a+b)\) ne se simplifie pas. Avant tout calcul logarithmique, vérifier que chaque argument est strictement positif.
\end{frame}
\begin{frame}{Cours de référence}
Pour résoudre \(a^x=b\), avec \(a>0\), \(a\neq1\) et \(b>0\), on prend le logarithme : \(x\log(a)=\log(b)\), donc \(x=\dfrac{\log(b)}{\log(a)}\).
\end{frame}
\begin{frame}{Cours de référence}
Résoudre \(1{,}08^x=2\) revient à écrire \(x=\dfrac{\log(2)}{\log(1{,}08)}\approx9{,}01\). Une valeur multipliée chaque année par \(1{,}08\) double donc après un peu plus de neuf ans.
\end{frame}
\begin{frame}{Définition du logarithme décimal}
\begin{itemize}
\item Relier puissance de 10 et logarithme.
\item Calculer des valeurs usuelles.
\item Utiliser la calculatrice.

\end{itemize}
\end{frame}
\begin{frame}{Définition du logarithme décimal}
Compléter les couples \(10^x=b\) et \(x=\log(b)\) pour \(b=0{,}001;0{,}1;1;10;1\,000\). Puis utiliser la calculatrice pour \(b=7\).
\end{frame}
\begin{frame}{Définition du logarithme décimal}
Pour \(b>0\), \(\log(b)\) est l'unique solution de \(10^x=b\). En particulier \(\log(1)=0\), \(\log(10)=1\) et \(\log(100)=2\).
\end{frame}
\begin{frame}{Définition du logarithme décimal}
Reprendre les données de l'activité et rédiger la solution dans l'ordre : définition des objets, propriété utilisée, calcul exact ou approché, puis phrase de réponse. La réponse modèle figure dans l'édition enseignant.
\end{frame}
\begin{frame}{Définition du logarithme décimal}
\textbf{Conduite en 1 h :} 7 min d'automatismes ; 10 min de mise en commun et trace ; 28 min d'exercices essentiels ; 5 min de correction ; 3 min de billet de sortie.

\textbf{Conduite en 2 h :} ajouter 20 min de recherche de l'activité, 20 min d'institutionnalisation, un second exercice autonome et 10 min de correction comparée.

\textbf{Différenciation :} fournir aux élèves fragiles une fiche à quatre cases (données, modèle, calcul, conclusion). Demander aux élèves rapides de produire un contre-exemple, une justification ou une seconde méthode.
\end{frame}
\begin{frame}{Sens de variation et représentation}
\begin{itemize}
\item Connaître le domaine.
\item Utiliser la croissance.
\item Lire la courbe.

\end{itemize}
\end{frame}
\begin{frame}{Sens de variation et représentation}
\begin{columns}[T,onlytextwidth]
\begin{column}{.57\textwidth}
Tracer \(y=10^x\) et la droite \(y=x\), puis obtenir la courbe du logarithme par symétrie. En déduire son domaine, son sens de variation et l'asymptote.
center center
\end{column}
\begin{column}{.4\textwidth}
\centering
\includegraphics[width=\linewidth,height=4.5cm,keepaspectratio]{figures/stmg-03-2.pdf}
\end{column}
\end{columns}
\end{frame}
\begin{frame}{Sens de variation et représentation}
La fonction \(\log\) est définie sur \(]0,+\infty[\), strictement croissante, et sa courbe admet \(x=0\) comme asymptote verticale.
\end{frame}
\begin{frame}{Sens de variation et représentation}
Reprendre les données de l'activité et rédiger la solution dans l'ordre : définition des objets, propriété utilisée, calcul exact ou approché, puis phrase de réponse. La réponse modèle figure dans l'édition enseignant.
\end{frame}
\begin{frame}{Sens de variation et représentation}
\textbf{Conduite en 1 h :} 7 min d'automatismes ; 10 min de mise en commun et trace ; 28 min d'exercices essentiels ; 5 min de correction ; 3 min de billet de sortie.

\textbf{Conduite en 2 h :} ajouter 20 min de recherche de l'activité, 20 min d'institutionnalisation, un second exercice autonome et 10 min de correction comparée.

\textbf{Différenciation :} fournir aux élèves fragiles une fiche à quatre cases (données, modèle, calcul, conclusion). Demander aux élèves rapides de produire un contre-exemple, une justification ou une seconde méthode.
\end{frame}
\begin{frame}{Propriétés algébriques}
\begin{itemize}
\item Développer ou condenser des logarithmes.
\item Vérifier les conditions.
\item Éviter les fausses propriétés.

\end{itemize}
\end{frame}
\begin{frame}{Propriétés algébriques}
À partir de \(10^2\times10^3=10^5\), conjecturer \(\log(ab)=\log a+\log b\). Tester numériquement avec \(a=2,b=5\), puis expliquer pourquoi \(a,b\) doivent être positifs.
\end{frame}
\begin{frame}{Propriétés algébriques}
Pour \(a,b>0\), \(\log(ab)=\log a+\log b\) et \(\log\!\left(\dfrac ab\right)=\log a-\log b\). De plus \(\log(a^n)=n\log a\).
\end{frame}
\begin{frame}{Propriétés algébriques}
Reprendre les données de l'activité et rédiger la solution dans l'ordre : définition des objets, propriété utilisée, calcul exact ou approché, puis phrase de réponse. La réponse modèle figure dans l'édition enseignant.
\end{frame}
\begin{frame}{Propriétés algébriques}
\textbf{Conduite en 1 h :} 7 min d'automatismes ; 10 min de mise en commun et trace ; 28 min d'exercices essentiels ; 5 min de correction ; 3 min de billet de sortie.

\textbf{Conduite en 2 h :} ajouter 20 min de recherche de l'activité, 20 min d'institutionnalisation, un second exercice autonome et 10 min de correction comparée.

\textbf{Différenciation :} fournir aux élèves fragiles une fiche à quatre cases (données, modèle, calcul, conclusion). Demander aux élèves rapides de produire un contre-exemple, une justification ou une seconde méthode.
\end{frame}
\begin{frame}{Résoudre \(10^x=b\)}
\begin{itemize}
\item Résoudre \(10^x=b\).
\item Encadrer une solution.
\item Interpréter un ordre de grandeur.

\end{itemize}
\end{frame}
\begin{frame}{Résoudre \(10^x=b\)}
Un appareil mesure des intensités \(b=0{,}04;3{,}2;850\). Pour chaque valeur, résoudre \(10^x=b\), donner \(x\) au centième et vérifier en calculant \(10^x\).
\end{frame}
\begin{frame}{Résoudre \(10^x=b\)}
L'équation \(10^x=b\) n'a de solution que si \(b>0\) ; cette solution est \(x=\log(b)\).
\end{frame}
\begin{frame}{Résoudre \(10^x=b\)}
Reprendre les données de l'activité et rédiger la solution dans l'ordre : définition des objets, propriété utilisée, calcul exact ou approché, puis phrase de réponse. La réponse modèle figure dans l'édition enseignant.
\end{frame}
\begin{frame}{Résoudre \(10^x=b\)}
\textbf{Conduite en 1 h :} 7 min d'automatismes ; 10 min de mise en commun et trace ; 28 min d'exercices essentiels ; 5 min de correction ; 3 min de billet de sortie.

\textbf{Conduite en 2 h :} ajouter 20 min de recherche de l'activité, 20 min d'institutionnalisation, un second exercice autonome et 10 min de correction comparée.

\textbf{Différenciation :} fournir aux élèves fragiles une fiche à quatre cases (données, modèle, calcul, conclusion). Demander aux élèves rapides de produire un contre-exemple, une justification ou une seconde méthode.
\end{frame}
\begin{frame}{Résoudre \(a^x=b\)}
\begin{itemize}
\item Résoudre \(a^x=b\).
\item Utiliser un quotient de logarithmes.
\item Arrondir avec sens.

\end{itemize}
\end{frame}
\begin{frame}{Résoudre \(a^x=b\)}
Le nombre d'abonnés est multiplié par \(1{,}15\) chaque mois. Déterminer au bout de combien de mois il aura triplé. Poser l'équation, la résoudre avec les logarithmes et interpréter l'arrondi.
\end{frame}
\begin{frame}{Résoudre \(a^x=b\)}
Pour \(a>0\), \(a\neq1\) et \(b>0\), \(a^x=b\) équivaut à \(x=\dfrac{\log(b)}{\log(a)}\).
\end{frame}
\begin{frame}{Résoudre \(a^x=b\)}
Reprendre les données de l'activité et rédiger la solution dans l'ordre : définition des objets, propriété utilisée, calcul exact ou approché, puis phrase de réponse. La réponse modèle figure dans l'édition enseignant.
\end{frame}
\begin{frame}{Résoudre \(a^x=b\)}
\textbf{Conduite en 1 h :} 7 min d'automatismes ; 10 min de mise en commun et trace ; 28 min d'exercices essentiels ; 5 min de correction ; 3 min de billet de sortie.

\textbf{Conduite en 2 h :} ajouter 20 min de recherche de l'activité, 20 min d'institutionnalisation, un second exercice autonome et 10 min de correction comparée.

\textbf{Différenciation :} fournir aux élèves fragiles une fiche à quatre cases (données, modèle, calcul, conclusion). Demander aux élèves rapides de produire un contre-exemple, une justification ou une seconde méthode.
\end{frame}
\begin{frame}{Équations \(x^a=b\) et inégalités simples}
\begin{itemize}
\item Résoudre une puissance inconnue.
\item Traiter une inégalité.
\item Tenir compte du domaine.

\end{itemize}
\end{frame}
\begin{frame}{Équations \(x^a=b\) et inégalités simples}
Résoudre sur \([0,+\infty[\) : \(x^3=40\), puis \(x^4=16\). Résoudre ensuite \(1{,}08^x<1{,}5\) et expliquer le rôle de la croissance.
\end{frame}
\begin{frame}{Équations \(x^a=b\) et inégalités simples}
Pour résoudre \(x^n=b\), tenir compte du domaine et de la parité de \(n\). Pour une inégalité exponentielle, le sens dépend de la croissance ou décroissance de la fonction.
\end{frame}
\begin{frame}{Équations \(x^a=b\) et inégalités simples}
Reprendre les données de l'activité et rédiger la solution dans l'ordre : définition des objets, propriété utilisée, calcul exact ou approché, puis phrase de réponse. La réponse modèle figure dans l'édition enseignant.
\end{frame}
\begin{frame}{Équations \(x^a=b\) et inégalités simples}
\textbf{Conduite en 1 h :} 7 min d'automatismes ; 10 min de mise en commun et trace ; 28 min d'exercices essentiels ; 5 min de correction ; 3 min de billet de sortie.

\textbf{Conduite en 2 h :} ajouter 20 min de recherche de l'activité, 20 min d'institutionnalisation, un second exercice autonome et 10 min de correction comparée.

\textbf{Différenciation :} fournir aux élèves fragiles une fiche à quatre cases (données, modèle, calcul, conclusion). Demander aux élèves rapides de produire un contre-exemple, une justification ou une seconde méthode.
\end{frame}
\begin{frame}{Problèmes contextualisés}
\begin{itemize}
\item Modéliser une situation.
\item Choisir une équation.
\item Conclure dans le contexte.

\end{itemize}
\end{frame}
\begin{frame}{Problèmes contextualisés}
Le niveau sonore est \(L=10\log\!\left(\dfrac I{I_0}\right)\). Calculer \(L\) lorsque \(I=10^6I_0\). Puis déterminer le rapport \(\dfrac I{I_0}\) pour \(L=85\) dB.
\end{frame}
\begin{frame}{Problèmes contextualisés}
Dans un problème, définir la variable, poser l'équation, résoudre, contrôler le domaine et conclure avec l'unité. Une solution numérique sans interprétation est incomplète.
\end{frame}
\begin{frame}{Problèmes contextualisés}
Reprendre les données de l'activité et rédiger la solution dans l'ordre : définition des objets, propriété utilisée, calcul exact ou approché, puis phrase de réponse. La réponse modèle figure dans l'édition enseignant.
\end{frame}
\begin{frame}{Problèmes contextualisés}
\textbf{Conduite en 1 h :} 7 min d'automatismes ; 10 min de mise en commun et trace ; 28 min d'exercices essentiels ; 5 min de correction ; 3 min de billet de sortie.

\textbf{Conduite en 2 h :} ajouter 20 min de recherche de l'activité, 20 min d'institutionnalisation, un second exercice autonome et 10 min de correction comparée.

\textbf{Différenciation :} fournir aux élèves fragiles une fiche à quatre cases (données, modèle, calcul, conclusion). Demander aux élèves rapides de produire un contre-exemple, une justification ou une seconde méthode.
\end{frame}
\begin{frame}{Bilan et évaluation}
\begin{itemize}
\item Synthétiser les méthodes.
\item Identifier les conditions.
\item Remédier aux erreurs.

\end{itemize}
\end{frame}
\begin{frame}{Bilan et évaluation}
Réaliser une carte de méthodes : domaine ; propriétés ; équations \(10^x=b\) ; équations \(a^x=b\) ; inégalités. Associer à chaque branche une condition et un exemple.
\end{frame}
\begin{frame}{Bilan et évaluation}
Avant un calcul de logarithme : vérifier la positivité. Avant une comparaison : identifier le sens de variation. Avant une conclusion : contrôler la vraisemblance.
\end{frame}
\begin{frame}{Bilan et évaluation}
Reprendre les données de l'activité et rédiger la solution dans l'ordre : définition des objets, propriété utilisée, calcul exact ou approché, puis phrase de réponse. La réponse modèle figure dans l'édition enseignant.
\end{frame}
\begin{frame}{Bilan et évaluation}
\textbf{Conduite en 1 h :} 7 min d'automatismes ; 10 min de mise en commun et trace ; 28 min d'exercices essentiels ; 5 min de correction ; 3 min de billet de sortie.

\textbf{Conduite en 2 h :} ajouter 20 min de recherche de l'activité, 20 min d'institutionnalisation, un second exercice autonome et 10 min de correction comparée.

\textbf{Différenciation :} fournir aux élèves fragiles une fiche à quatre cases (données, modèle, calcul, conclusion). Demander aux élèves rapides de produire un contre-exemple, une justification ou une seconde méthode.
\end{frame}
\end{document}
